%%%%%
%
\section{Correction d'un param\`etre \texttt{Int} par le rapport des masses volumiques
$\rho / \rho_0$ }

\itemstyle{Op\'erateur~:} {\sf Cor\_dens\_int}
\index{Cor\_dens\_int}

%%%%%OBJET%%%%
\itemstyle{Objet~:}
Cet op�rateur multiplie un param\`etre de type \texttt{Int} par le facteur $\rho/\rho_0$, qui est le
rapport entre la masse volumique de l'air ambiant et la masse volumique de
l'air sec standard \`a 1013.5 hPa et 15 $^{\circ}C$.   

%%%%%ARGUMENTS%%%%
\bigskip
%%%%%ENTREES%%%%
\begin{tabarg}
\multicolumn{3}{l}{\titretab {Entr\'ees}}\\
\hline
$Var_{brut}$ & Vecteur, {\tt int16} & Variable \`a corriger (unit\'e quelconque)\\
$Dim_{Var_{brut}}$ & Coeff., {\tt int} & Dimension de la variable \`a corriger.\\
$\rho$ & Vecteur, {\tt float} & densit\'e ($kg/m^3$)\\
\end{tabarg}

\medskip
%%%%%SORTIES%%%%
\begin{tabarg}
\multicolumn{3}{l}{\titretab {Sorties}}\\
\hline
$Var_{cor}$ & Vecteur, {\tt float} & Variable corrig\'ee (m�me unit\'e que $Var_{brut}$)\\
\end{tabarg}


%%%%%FORMULE%%%%
\itemstyle{Formule}

%Cor_dens
\[
Var_{cor} = \frac{\rho}{\rho_0} * Var_{brut}
\]

avec $\rho_0 = 1.22500\ kg.m^{-3}$ (par d\'efinition).


%%%%%COMMENTAIRE%%%%
\itemstyle{Commentaire :}
Cette correction est utilis\'ee pour certaines sondes
micro-physique (PCASP, TSI) lorsqu'une la pompe de r�gulation du d�bit 
est asservie en "masse constante" 
alors qu'on souhaite travailler \`a "volume constant". 

%%%%%AUTEUR%%%%
\itemstyle{Auteur~:}
CNRM/GMEI/TRAMM

%%%%%REFERENCE%%%%
\itemstyle{R\'ef\'erences~:}
Principe de la correction :\\
Fichier ADJUSCF.TXT du 04/12/1997, par Ch. Allet (INSU/Voloth�que).\\
Valeur de $\rho_0$ : \\
Standard ISO2533-1975 (Atmosph�re-type) \cite{ISO-2533}, tableau 5.
